\subsection{Selección de Características}

Utiliza alguno de los conuntos de datos propuestos:

\textbf{Arritmia.}  ($Arrhythmia$ $Database$ $UCI$ $Machine$ $Learning$ $Repository$), el objetivo de este conjunto de datos es
diferenciar entre arritmias cardíacas y clasificarlas en uno de los 16 grupos. La clase 01 agrupa a las señales de electrocardiograma
(ECG) 'normales', del 02 al 15 a diferentes clases de arritmias y la clase 16 al resto de las señales ECG de arritmias no clasificadas.
Cuenta con 452 muestras, 279 características, 206 de las cuales son de valores cuantitativos y el resto cualitativos.

\textbf{Expresión genética en pacientes con leucemia} ($Gene$ $expression$ $dataset$ $from$ $Golub$ $et$ $al.$ $(1999)$),
este conjunto de datos consta de 47 pacientes con leucemia linfoblástica aguda (LLA) y 25 pacientes con leucemia mieloide aguda
(LMA). A cada uno de los 72 pacientes se le extrajeron muestras de médula ósea al momento del diagnóstico, las mediciones de expresión
génica de las muestras resultaron en 7129 características que representan los niveles de expresión de genes para cada muestra. Este
conjunto de datos tiene una alta dimensión con pocas muestras, valores con ruido y posible desbalance de clases.

\begin{enumerate}
    \item Verificar la integridad del conjunto de datos desde valores faltantes, clases minoritarias con muy pocas muestras, datos
          atípicos y escalas de datos diferentes (algunas características son frecuencias, otras voltajes).

    \item{Preparación de Datos}
          \begin{enumerate}
              \item Detectar valores faltantes y completarlos si es el caso.
              \item  Aplicar escaladores para remover o mitigar el efecto de muestras atípicas (por ejemplo z-score).
              \item  Aplicar técnicas para corregir desbalanceos (como sobremuestreo de la clase minoritaria).
              \item Implementar validación cruzada para evaluar cada método de selección manteniendo un conjunto de prueba separado.
              \item  Dividir los datos (80\% entrenamiento / 20\% prueba) sin incluir los datos de prueba en la optimización de parámetros.
          \end{enumerate}

    \item{Implementación de Métodos}
          \begin{enumerate}
              \item \textbf{Búsqueda Secuencial hacia delante (SFS)}
                    \begin{itemize}
                        \item SFS con validación cruzada e interrupción temprana ($early$ $stopping$) si la exactitud no mejora en 3 iteraciones.
                        \item Generar gráfico exactitud por características (20 mejores características).
                    \end{itemize}

              \item \textbf{Fisher}
                    \begin{itemize}
                        \item Implementar función del cálculo  Score de separabilidad entre clases.
                        \item Calcular puntuaciones de relevancia de cada característica.
                        \item Selecciona las 20 características con mayor puntuación.
                    \end{itemize}

              \item \textbf{PCA}
                    \begin{itemize}
                        \item Remover pares de características altamente correlacionadas (coeficiente $>$ 0.9).
                        \item Determinar componentes para 95\% varianza.
                        \item Visualizar proyección 2D.
                    \end{itemize}
          \end{enumerate}

    \item{Análisis de resultados}
          \begin{enumerate}
              \item ¿Qué método obtuvo mejor exactitud? ¿Por qué?.

              \item ¿Qué método es más robusto ante outliers y desbalanceo?

              \item ¿Cómo afecta la correlación entre características al PCA?.

              \item ¿Qué ventajas tiene Fisher Score sobre SFS en términos computacionales?.
          \end{enumerate}
\end{enumerate}


